\exo

Calculer la moyenne de la série statistique suivante.

   \begin{tabular}{| l | c | c | c | c | c | c |}
     \hline
     Valeurs $x_i$ & 17 & 20 & 25 & 35 & 48 & 60 \\ \hline
     Effectifs $n_i$ & 5 & 10 & 5 & 8 & 12 & 13 \\ 
     \hline
   \end{tabular}
\exo

Calculer la moyenne de la série statistique suivante.

   \begin{tabular}{| l | c | c | c | c | c |}
     \hline
     Valeurs $x_i$ & 8 & 9 & 12 & 17 & 20 \\ \hline
     Fréquences $f_i$ & 0,13 & 0,26 & 0,08 & 0,16 & 0,37 \\ 
     \hline
    \end{tabular}
    
\exo
L'entraîneur d'un club de football a publié les statistiques des matchs de la saison.
Sans calculatrice, calculer le nombre moyen de buts marqués par match.

   \begin{tabular}{| l | c | c | c | c |}
     \hline
     Nombre de buts par match & 0 & 1 & 2 & 3 \\ \hline
     Fréquences & 0,35 & 0,15 & 0,4 & 0,1\\ 
     \hline
    \end{tabular}

\exo

 \begin{tabular}{| l | c | c | c | c |}
     \hline
     Valeurs $x_i$ & $x_1$ & $x_2$ & $x_3$ & $x_4$ \\ \hline
     Effectifs $n_i$ & $n_1$ & $n_2$ & $n_3$ & $n_4$ \\ 
     \hline
    \end{tabular}

\begin{enumerate}
\item Donner l'expression de la moyenne $m$ de cette série.
\item On multiplie toutes les valeurs par un réel $k$. Présenter la nouvelle série sous la forme d'un tableau.
\item Exprimer la moyenne $m'$ de cette série.
\item En factorisant par $k$, justifier que $m'=k\times m$.
\end{enumerate}

\exo

Dans une très petite entreprise, les salaires mensuels des employés, en euros, sont les suivants : 1173 ; 1173 ; 1190 ; 1205 ; 1235 ; 3200.

\begin{enumerate}
\item Calculer la moyenne des salaires dans l'entreprise.
\item Le salaire le plus élevé passe à $4850$ euros. Calculer la nouvelle moyenne.
\item Expliquer la phrase : \og La moyenne est sensible aux valeurs extrêmes. \fg{}
\end{enumerate}


\exo
Calculer la variance et l'écart-type de la série suivante :
\[
57\ ;\ 59\ ;\ 62\ ;\ 87\ ;\ 98\ ;\ 100.
\]
\exo
Calculer la variance et l'écart-type de la série suivante.

\begin{tabular}{| l | c | c | c | c |}
     \hline
     Valeurs $x_i$ & 10 & 12 & 17 & 20 \\ \hline
     Effectifs $n_i$ & 1 & 4 & 3 & 2  \\ 
     \hline
   \end{tabular}
\exo
On a relevé les tailles, en centimètres, des membres d'un groupe de danse folklorique.
   
   \begin{tabular}{| l | c | c | c | c | c | c |}
     \hline
     Taille (en cm) & 148 & 152 & 155 & 160 & 162 & 164 \\ \hline
     Effectifs & 1 & 3 & 5 & 2 & 5 & 3  \\ 
     \hline
   \end{tabular}
   
   \begin{enumerate}
\item Calculer la taille moyenne des membres du groupe.
\item Calculer l'écart-type de cette série.
   \end{enumerate}
   
\exo
   Pour chacune des séries suivantes, déterminer la médiane et les quartiles.
   \begin{enumerate}
   \item 140 ; 140 ; 148 ; 149 ; 150 ; 152 ; 152 ; 160 ; 170 ; 170 ; 175 ; 180
   \item 14 ; 15 ; 15 ; 16 ; 17 ; 17 ; 17 ; 18 ; 20 ; 20 ; 20 ; 20 ; 25
   \item \begin{tabular}{| l | c | c | c | c | c | c |}
     \hline
     Valeurs $x_i$ & 6 & 7 & 9 & 11 & 12 & 16 \\ \hline
     Effectifs $n_i$ & 2 & 3 & 4 & 8 & 5 & 3  \\ 
     \hline
   \end{tabular}
   \end{enumerate}
   
\exo
   \begin{enumerate}
   \item Construire une série de neuf valeurs telle que $Q_1=3$, $Me=8$ et $Q_3=10$.
   \item Construire une série de cinq valeurs de moyenne 40, de médiane 20 et d'étendue 100. 
   \end{enumerate}
   
\exo
   On considère la série suivante :
   \[
   18\ ;\ 20\ ;\ 20\ ;\ 22\ ;\ 24\ ;\ 24\ ;\ 28\ ;\ 30\ ;\ 30\ ;\ 30.
   \]
   Que deviennent la médiane et les quartiles de cette série :
   \begin{enumerate}
   \item si l'on ajoute 10 à toutes les valeurs ?
   \item si l'on multiplie toutes les valeurs par 10 ?
   \end{enumerate}
   
\exo
   On donne les températures moyennes mensuelles, en degrés Celsius, de deux villes.

\begin{center}
\small
\renewcommand{\arraystretch}{0.95}
\setlength{\tabcolsep}{4pt}
\begin{tabular}{|l|c|c||l|c|c|}
     \hline
Mois & Mexico & Barcelone & Mois & Mexico & Barcelone \\
     \hline
Janvier & 12,4 & 9,5 & Juillet & 16,7 & 24,5 \\ \hline
Février & 14,1 & 10,3 & Août & 16,8 & 24,3 \\ \hline
Mars & 16,2 & 12,4 & Septembre & 16,3 & 21,8 \\ \hline
Avril & 17,4 & 14,6 & Octobre & 15,1 & 17,6 \\ \hline
Mai & 18,4 & 17,7 & Novembre & 13,9 & 13,5 \\ \hline
Juin & 17,7 & 21,5 & Décembre & 12 & 10,3 \\
     \hline
\end{tabular}
\end{center}
   
     Pour la ville de Mexico, calculer :
  
{\small
  \begin{enumerate}[itemsep=0pt, topsep=2pt, parsep=0pt]

  \item l'étendue des températures ;
  \item la température annuelle moyenne ;
  \item la médiane et les quartiles $Q_1$ et $Q_3$.
\item Pour Barcelone, on donne : moyenne $16,48$ degrés Celsius, étendue $14,8$ degrés Celsius, médiane $16,1$ degrés Celsius, $Q_1=10,3$ et $Q_3=21,5$.
À l'aide de ces indicateurs, répondre aux questions suivantes.
  \item il fait plus chaud à Barcelone qu'à Mexico ?
  \item les écarts de température sont-ils moindres à Mexico ?
  \item dans ces deux villes, la température est-elle supérieure à 16 degrés Celsius pendant au moins la moitié de l'année ?
  \item à Mexico, fait-il environ entre 14 et 17 degrés Celsius pendant la moitié de l'année ?
  
  \end{enumerate}
}
